\documentclass{homework}
% \author{Tashfeen, Ahmad}  % uncomment with your name
\class{CSCI 3223: Tashfeen's Machine Learning II}
\title{Homework 5}
\address{
  Oklahoma City University,
  Petree College of Arts \& Sciences,
  Computer Science
}

\usepackage{blkarray}
\DeclareMathOperator*{\argmax}{argmax}

\begin{document} \maketitle

You're \textit{not allowed} to use any libraries other than Numpy or
Matplotlib.

\question Read Mitchell, Chapter 13 (pg. 267--377). What do you think
sets \textit{Reinforcement Learning} apart from other machine learning
techniques? Do you see any pitfalls? Give at least one.


\question\label{toy} Consider a toy guessing-game-environment. The environment
keeps track of two integer quantities $0< (x, \hat{x}) < 99$. The
agent's aim is to guess the number $x$ and the agent's guess is
$\hat{x}$. We define the state space as,
\[
  S = \{\hat{x} < x, \hat{x} = x, \hat{x} > x\}.
\]
And, similarly, the action space,
\[
  A = \curl{
    \hat{x} \leftarrow \hat{x} + \floor{\frac{\hat{x}}{2}},
    \hat{x} \leftarrow \hat{x},
    \hat{x} \leftarrow \hat{x} - \floor{\frac{\hat{x}}{2}}}.
\]
Upon each action, the environment returns the following reward,
\[
  r =
  \begin{cases}
    +2, \text{ if } \hat{x} = x, \\
    +1, \text{ if } \hat{x} \text{ is closer to } x \text{ than before,}\\ 
    -1, \text{ if } \hat{x} \text{ is farther than before or equal to } x
  \end{cases}
\]

Using the learning rule given in Equation 13.7,
\[
  \hat{Q}(s, a) = r + \gamma \max_{a'}\hat{Q}(s', a')
\]
and the algorithm in Table 13.1 on Mitchell pg. 375, give the
converged $\hat{Q}$ table.  Use the initial values,
\[
  \hat{x} = 64, \quad x = 4, \quad \gamma = 0.5, \quad \hat{Q}(s, a) = 0 \text{ for all } (s, a)
\]
and the following policy function,
\[
  \pi(s) = \argmax_a\hat{Q}(s, a).
\]
Place the code for this question in \texttt{toy.py}.

\question Does question \ref{toy} describe a Markov Decision Process (MDP)?
Briefly explain your answer. Hint: Mitchell pg. 370.

\question Assume that $\hat{Q}$ converges to,
\[
  \begin{blockarray}{cccc}
    \uparrow & \nil & \downarrow & {}\\
    \begin{block}{[ccc]c}
      2    & 0  & 0 & \hat{x} < x \\
      -0.5 & 4  & 0 & \hat{x} = x \\
      -1   & -1 & 2 & \hat{x} > x \\
    \end{block}
  \end{blockarray}
\]
Interpret the policy $\pi(s) = \argmax_a\hat{Q}(s, a)$. Hint: If the
$\hat{Q}$ from your question \ref{toy} did not converge to the one
given in this question, you have made a mistake.

\question Write a $Q$-learning snake agent for the following game environment,\\
\url{https://github.com/simurgh9/snake}\\
You should read the documentation \texttt{README.md} to get started.

Note that this question will be graded based on your code's quality
and agent's performance. Please avoid duplicating code. Ask the
professor if you need some logic that you think is missing.

Give the details of hyper parameters and your state \& action state
modelling. Give any other details necessary.

\section*{Submission Instructions}
\begin{enumerate}
  \item Submit a PDF that answers all the questions including any plots.
  \item Submit the contents of the director that you created under
  \texttt{snake/model/serpents/}. Also submit the Python file
  \texttt{toy.py} from question \ref{toy}.
  \item Submit an optimised \texttt{demo.gif} of your $Q$-learning
  agent playing the game. An example of such a gif is also in the
  \texttt{README.md}
\end{enumerate}
\end{document}

